%==============================================================================
% Preambulo — Tesis de Grado UTN FRM
% Formato: APA 7ma edicion adaptado a tesis
% Compilador: XeLaTeX
%==============================================================================

% --- Clase de documento ---
\documentclass[11pt,a4paper]{article}

% --- Idioma: Espanol ---
\usepackage[spanish,es-tabla,es-nodecimaldot,es-noindentfirst]{babel}

% --- Fuentes: Arial 11pt ---
\usepackage{fontspec}
\setmainfont{Arial}
\setsansfont{Arial}
\setmonofont{Courier New}

% --- Geometria: 2.54 cm (1 pulgada) todos los margenes (APA 7) ---
\usepackage[paper=a4paper,left=2.54cm,right=2.54cm,top=2.54cm,bottom=2.54cm]{geometry}

% --- Interlineado: Doble espacio (APA 7) ---
\usepackage{setspace}
\doublespacing

% --- Sangria en primera linea de parrafos: 1.27 cm (APA 7) ---
\usepackage{indentfirst}
\setlength{\parindent}{1.27cm}
\setlength{\parskip}{0pt}

% --- Encabezados y pies de pagina ---
\usepackage{fancyhdr}
\pagestyle{fancy}
\fancyhf{}
\fancyhead[L]{\small AUTOMATIZACION DE MESA DE AYUDA}
\fancyhead[R]{\small\thepage}
\renewcommand{\headrulewidth}{0.4pt}
% Primera pagina de cada seccion sin encabezado (APA 7 permite running head o no)
\fancypagestyle{plain}{%
  \fancyhf{}
  \fancyhead[R]{\small\thepage}
  \renewcommand{\headrulewidth}{0pt}
}

% --- Titulos de seccion: APA 7 niveles ---
\usepackage{titlesec}
% Nivel 1: Centrado, negrita, titulo en mayuscula/minuscula
\titleformat{\section}[block]{\centering\bfseries\normalsize}{\thesection.}{0.5em}{}
\titlespacing*{\section}{0pt}{12pt}{6pt}
% Nivel 2: Alineado a izquierda, negrita
\titleformat{\subsection}[hang]{\bfseries\normalsize}{\thesubsection.}{0.5em}{}
\titlespacing*{\subsection}{0pt}{12pt}{6pt}
% Nivel 3: Alineado a izquierda, negrita cursiva
\titleformat{\subsubsection}[hang]{\bfseries\itshape\normalsize}{\thesubsubsection.}{0.5em}{}
\titlespacing*{\subsubsection}{0pt}{12pt}{6pt}

\usepackage{tocloft}
\renewcommand{\cftsecleader}{\cftdotfill{\cftdotsep}}
\renewcommand{\cftdotsep}{1}

% --- Tablas: formato APA 7 (solo bordes horizontales, sin lineas verticales) ---
\usepackage{booktabs}
\usepackage{array}
\usepackage{multirow}
\renewcommand{\arraystretch}{1.2}

% --- Figuras ---
\usepackage{graphicx}
\graphicspath{{./figures/}}
\usepackage[font=small,labelfont=bf,labelsep=period,justification=raggedright,singlelinecheck=false]{caption}

% --- Ecuaciones matematicas y simbolos ---
\usepackage{amsmath,amssymb,amsfonts}
\usepackage{siunitx}
\sisetup{output-decimal-marker={,},group-separator={.},group-minimum-digits=4}

% --- Comillas tipograficas (requerido por biblatex) ---
\usepackage[autostyle,spanish=mexican]{csquotes}

% --- Hipervinculos (APA 7: sin bordes de color) ---
\usepackage[hidelinks,bookmarks=true,bookmarksnumbered=true,pdfstartview=FitH]{hyperref}
\hypersetup{
  pdftitle={Automatizacion inteligente del registro de incidentes en mesas de ayuda},
  pdfauthor={Luca Gomez, Carla Bustos, Gonzalo Sevilla},
  pdfsubject={Tesis de Grado — Tecnicatura Universitaria en Programacion},
  pdfkeywords={automatizacion, mesa de ayuda, NLP, LLM, N8N, FastAPI, clasificacion de tickets, espanol rioplatense},
  colorlinks=false,
  pdfborder={0 0 0}
}

% --- Bibliografia: APA 7ma edicion ---
\usepackage[style=apa,backend=biber,sorting=nyt,language=spanish]{biblatex}
\addbibresource{../Bibliography_base.bib}
% Mapeo de campos para APA 7 en espanol
\DefineBibliographyStrings{spanish}{%
  bibliography = {Referencias bibliograficas},
  references = {Referencias},
  page = {p.},
  pages = {pp.},
  andothers = {et al.},
  retrieved = {Recuperado},
  from = {de},
  urlseen = {Consultado},
  edition = {ed.},
  volume = {Vol.},
  number = {N.°},
  in = {En},
}

% --- Listas enumeradas ---
\usepackage{enumitem}
\setlist{nosep,labelindent=1.27cm,leftmargin=*}

% --- Codigo fuente ---
\usepackage{listings}
\lstset{
  basicstyle=\ttfamily\small,
  breaklines=true,
  frame=single,
  numbers=left,
  numberstyle=\tiny,
  showstringspaces=false,
  tabsize=2,
  captionpos=b
}

% --- Apendices ---
\usepackage[toc,title]{appendix}

% --- Comandos personalizados ---
\newcommand{\doi}[1]{\href{https://doi.org/#1}{#1}}
\newcommand{\urlref}[1]{\href{#1}{#1}}

% --- Metadatos del documento ---
\title{Automatizacion inteligente del registro de incidentes\\en mesas de ayuda empresariales mediante\\orquestacion de flujos con N8N y procesamiento\\de lenguaje natural basado en modelos\\de lenguaje grandes}
\author{}
\date{2026}
